% Part: set-theory
% Chapter: ordinals
% Section: introduction

\documentclass[../../../include/open-logic-section]{subfiles}

\begin{document}

\olfileid{sth}{ordinals}{intro}

\olsection{प्रस्तावना}

\olref[z][]{chap} मध्ये \stagesinf{} या रूपात पदानुक्रमात अनंत टप्पा आहे असे आपण गृहीत धरले (अनंततेचे स्वयंसिद्धकही पाहा). पण \stagessucc{} मानल्यावर त्या टप्प्यावर थांबता येत नाही; पुढे जावेच लागते. म्हणून पहिल्या अनंत टप्प्यानंतरच्या टप्प्यावर त्या पहिल्या अनंत टप्प्यावर उपलब्ध असलेल्या संचांपासून बनणारे सर्व शक्य संग्रह तयार होतात; मग पुन्हा तेच; पुन्हा; पुन्हा; \dots

यात अप्रकटपणे आपण संख्येची एक ``सहज'' संकल्पना वापरू लागलो आहोत: सर्व नैसर्गिक संख्यांच्या \emph{नंतरही} संख्या असू शकतात. विशेषतः येथे अभिप्रेत आहे ती \emph{अनंतपार क्रमसूचक संख्या}. या कल्पनेला अधिक काटेकोर रूप देणे हे या प्रकरणाचे उद्दिष्ट आहे. आपण क्रमसूचक संख्यांची सर्वसाधारण संकल्पना पाहू आणि मग काही विशिष्ट संचांची स्पष्ट व्याख्या आपल्या क्रमसूचक संख्या म्हणून करू.

\end{document}
